\chapter{Gode råd}

\section{\latex  kan meget mere end du tror!}

Mange nybegyndere tror, at de skal bruge Word-lignende krumspring for
at få et bestemt udseende. Hvis det ikke lykkes for dem, konkluderer
de ofte, at \latex må lide af alvorlige mangler.

\subsection{Typiske misforståelser}

\begin{itemize}
\item Måske er det layout, man ønsker sig, ikke hensigtsmæssigt. Det
  er almindeligt at man har vaner med sig fra gymnasiet og fra
  Word-land, som snarere skyldes uvaner end stor typografisk
  indsigt. Det er langt bedre at tage ved lære af typografien i
  fagbøger og videnskabelige publikationer end at lade sig styre af en
  gammel SRP-opgave.
  
\item Og der er med meget stor sandsynlighed snarere tale om at der er
  noget, man ikke ved selv, end om noget, det er umuligt at gøre ved
  brug af \latex.
\end{itemize}

\subsection{Noget om layout}

Hvis man \emph{slet ikke} kan acceptere standard-layoutet og drømmer sig
tilbage til SRP-dagene, er der pakken \verb2wordlike2. Her er stor
linjeafstand, Times New Roman, smalle marginer og andet godt, så man
kan genskabe det trygge gymnasielayout og drømme om Word-pølse.

\begin{ovelse}
  Prøv at skrive \verb2\usepackage{wordlike}2 i præamblen, og se hvad
  der sker.
\end{ovelse}

Mange nybegyndere tror også, at de er nødt til at lave manuelle
krumspring mange steder i et dokument for at styre dets layout. Typisk
ved de ikke, at \verb2report2-dokumentklassen findes og kender kun
\verb2article2. Derfor bruger de tvungne sideskift med
\verb2\newpage2. Nybegyndere synes også ofte, at der skal være større
afstand mellem paragraffer og indsætter tvungne linjeskift med
\verb2\\2 overalt.

Dette er en anden typisk arv fra Word-land, hvor man er nødt til at
skulle håndtere den slags manuelt rundt om i Word-pølsen. Men i \latex
er der en masse parametre, der bestemmer marginbredde, linjeafstand
mm. Dem kan man selv sætte én gang for alle oppe i præamblen, så man
ved en lille ændring dér kan ændre dokumentets layout fuldstændig.

Se også
\url{https://da.overleaf.com/learn/latex/Articles/How_to_change_paragraph_spacing_in_LaTeX},
der forklarer de forskellige parametre, og hvordan man sætter dem ét
sted i dokumentet.

Her er nogle små øvelser, der viser hvad parametrene kan.

\begin{ovelse}
  Prøv at skrive \verb2\documentclass[a4paper,11pt,twocolumn]{report}2 i præamblen og se,
  hvad der sker.
\end{ovelse}

\begin{ovelse}
  Prøv at skrive
\begin{verbatim}
\setlength{\parskip}{4mm}
\end{verbatim}
  i præamblen og se, hvad der sker.
\end{ovelse}

\begin{ovelse}
  Prøv at skrive
\begin{verbatim}
\setlength{\parindent}{0mm}
\end{verbatim}
  i præamblen og se, hvad der sker.
\end{ovelse}

\begin{ovelse}
  Prøv at skrive
  \begin{verbatim}
\renewcommand\baselinestretch{1.5}
\end{verbatim}
  i præamblen og se, hvad der sker.
\end{ovelse}

%%% Local Variables:
%%% mode: latex
%%% TeX-master: "report"
%%% End:
